% !TeX root = ../navier-stokes-manuscript.tex
\subsection{The first rectangle and the exact unpaid integral}\label{sub:BodyCompatibleIntersection}

The trace argument has reduced an infinite smoothness question to a finite
place where the equation can be tested.  At \((N,M)=(0,1)\), the velocity must
sit one spatial step above \(H^1\), and its first time response must sit one
step below:
\begin{equation}\label{eq:BodyFirstWholeTerminalRectangle}
  u\in L^2(I_*;H^2),
  \qquad
  u_t\in L^2(I_*;L^2).
\end{equation}
Together these two rows reach \(u(T_*)\) in \(H^1\).  We now ask the original
equation what quantity has to remain finite for this lowest rectangle to reach
the terminal face.

Write
\begin{equation}\label{eq:BodyEnstrophyDissipationQuantities}
\begin{aligned}
  Z(t)&:=\norm{\Lambda u(t)}_2^2=\norm{\omega(t)}_2^2,
  &
  Y(t)&:=\norm{\Lambda^2u(t)}_2^2=\norm{\nabla\omega(t)}_2^2,
  \\
  D(t)&:=\norm{\Lambda^{3/2}u(t)}_2^2.
\end{aligned}
\end{equation}
The enstrophy balance already contains the spatial row \(Y\); the only term
capable of preventing its integration is vortex stretching.  Since the strain
is recovered from vorticity by Riesz transforms, H\"older and the Sobolev
embeddings give
\begin{align}
  \left|\int_{\R^3}\omega\cdot S\omega\,dx\right|
  &\leq
  \norm{\omega}_2\norm{\omega}_3\norm{S}_6
  \notag\\
  &\leq C Z^{1/2}D^{1/2}Y^{1/2}.
  \label{eq:BodyFirstRectangleStretchingEstimate}
\end{align}
Young's inequality returns half of \(Y\) to viscosity and leaves the critical
coefficient beside \(Z\):
\begin{equation}\label{eq:BodyEnstrophyFromCriticalDissipation}
  Z'(t)+\nu Y(t)
  \leq C\nu^{-1}D(t)Z(t).
\end{equation}
The coefficient has isolated the endpoint interface:
\begin{equation}\label{eq:BodyCriticalDissipationInterface}
  D_*:=\int_0^{T_*}D(t)\,dt
  =\int_0^{T_*}\norm{\Lambda^{3/2}u(t)}_2^2\,dt.
\end{equation}

Suppose first that \(D_*<\infty\).  Gronwall applied to
\eqref{eq:BodyEnstrophyFromCriticalDissipation} gives
\begin{equation}\label{eq:BodyEnstrophyAndHtwoFromDissipation}
  \sup_{0<t<T_*}Z(t)+\int_0^{T_*}Y(t)\,dt<\infty.
\end{equation}
Together with the kinetic-energy row, this places
\(u\in L^2(I_*;H^2)\).  The projected momentum equation reads
\[
  u_t=\nu\Delta u-\Leray\bigl((u\cdot\nabla)u\bigr),
\]
and its nonlinear part satisfies
\begin{equation}\label{eq:BodyNonlinearityFromCriticalDissipation}
  \norm{\Leray((u\cdot\nabla)u)}_2^2
  \leq C\norm{u}_6^2\norm{\nabla u}_3^2
  \leq C ZD.
\end{equation}
The uniform \(Z\) roof and the integrable \(D\) coefficient place
\(u_t\in L^2(I_*;L^2)\).  This is exactly
\eqref{eq:BodyFirstWholeTerminalRectangle}.

Conversely, a cutoff-uniform \((0,1)\) rectangle places \(u\) in
\(L^2(I_*;H^2)\) by monotone convergence, and the intermediate
\(H^{3/2}\) row is then integrable.  The lower Sobolev pieces are already paid
by the kinetic-energy identity in
\Cref{prop:ClassicalKineticEnergyIdentity}.  We have proved the exact
equivalence
\begin{equation}\label{eq:BodyFirstRectangleCriticalDissipationEquivalence}
  D_*<\infty
  \quad\Longleftrightarrow\quad
  \sup_{0<T<T_*}\mathfrak R_{0,1}(u;T)<\infty.
\end{equation}

The same interface can be read directly in the critical-height balance.
Because \(D=2E_{3/2}\), its time integral is
\begin{equation}\label{eq:BodyCriticalPrefixIdentity}
  H(t)+\nu\int_0^tD(s)\,ds
  =H(0)+\int_0^t\mathcal P_{1/2}(s)\,ds.
\end{equation}
An upper bound on the signed production prefix makes \(D_*\) finite because
\(H\geq0\).  In the other direction, \(D_*<\infty\) gives the first rectangle,
the adjacent trace in \Cref{lem:AdjacentHilbertScaleTrace} gives a bounded
\(H\), and the same equality bounds the prefix.  Thus three descriptions meet
at one interface:
\begin{equation}\label{eq:BodyFirstRectangleSignedPrefixEquivalence}
  \sup_{0<t<T_*}\int_0^t\mathcal P_{1/2}(s)\,ds<\infty
  \quad\Longleftrightarrow\quad
  D_*<\infty
  \quad\Longleftrightarrow\quad
  \sup_{0<T<T_*}\mathfrak R_{0,1}(u;T)<\infty.
\end{equation}
The signed prefix is not a cheaper route around the first rectangle.  It is
the same missing payment written in the coordinates of the vortex.

There is a regime in which the datum makes that payment.  Set
\[
  a(t):=\norm{\Lambda^{1/2}u(t)}_2,
  \qquad
  d(t):=\norm{\Lambda^{3/2}u(t)}_2^2.
\]
For \(\mathcal B(u,u):=\Leray((u\cdot\nabla)u)\), the dual Sobolev embeddings
give
\begin{align}
  \norm{\mathcal B(u,u)}_{\dot H^{-1/2}}
  &\leq C\norm{u\cdot\nabla u}_{3/2}
  \notag\\
  &\leq C\norm{u}_3\norm{\nabla u}_3
  \leq C a(t)d(t)^{1/2}.
  \label{eq:BodyCriticalNonlinearDualEstimate}
\end{align}
Pairing this source with the remaining half derivative of the velocity gives
\begin{equation}\label{eq:BodyCriticalProductionAbsoluteEstimate}
  |\mathcal P_{1/2}(t)|
  \leq C a(t)d(t).
\end{equation}
The critical-height row becomes
\begin{equation}\label{eq:BodySmallCriticalDataAbsorption}
  \frac12\frac d{dt}a(t)^2
  +\bigl(\nu-Ca(t)\bigr)d(t)
  \leq0.
\end{equation}
If \(a(0)<\nu/C\), continuity prevents a first crossing of the viscous
threshold.  Before such a crossing, the coefficient is positive and \(a\) is
nonincreasing, so the crossing cannot occur.  Hence
\begin{equation}\label{eq:BodySmallCriticalDataDissipation}
  \frac12a(t)^2
  +\bigl(\nu-Ca(0)\bigr)\int_0^t d(s)\,ds
  \leq\frac12a(0)^2.
\end{equation}
Small critical data pay \(D_*\) from the beginning.  An arbitrary Schwartz
datum need not begin below this threshold, and scaling cannot improve the
coefficient: nonlinear production and critical dissipation change with the
same weight.

Kinetic energy can still carry a long-lived classical history into the small
regime.  This produces a horizon beyond which a finite maximal time cannot
sit.
\begin{corollary}[Small-data capture horizon]
\label[corollary]{cor:BodySmallDataCaptureHorizon}
Let \(C\) be the constant in
\eqref{eq:BodySmallCriticalDataAbsorption}.  If a nonzero datum
\(u_0\in\Ssigma\) generates a finite maximal time, then
\begin{equation}\label{eq:BodySmallDataCaptureHorizon}
  T_*
  \leq
  \frac{C^4\norm{u_0}_2^4}{2\nu^5}.
\end{equation}
\end{corollary}

\begin{proof}
Fix \(0<T<T_*\).  The kinetic-energy identity in
\Cref{prop:ClassicalKineticEnergyIdentity} gives
\[
  \int_0^T Z(t)\,dt
  \leq\frac{\norm{u_0}_2^2}{2\nu}.
\]
At some \(t_0\in(0,T)\),
\[
  Z(t_0)\leq\frac{\norm{u_0}_2^2}{2\nu T}.
\]
Interpolation between the energy and enstrophy heights gives
\[
  a(t_0)^2
  \leq\norm{u(t_0)}_2\norm{\Lambda u(t_0)}_2,
  \qquad
  a(t_0)
  \leq\frac{\norm{u_0}_2}{(2\nu T)^{1/4}}.
\]
If \(T>C^4\norm{u_0}_2^4/(2\nu^5)\), the history at \(t_0\) lies below
\(\nu/C\).  Restart the preceding first-crossing argument there:
\[
  \frac12a(t)^2
  +\bigl(\nu-Ca(t_0)\bigr)\int_{t_0}^t d(s)\,ds
  \leq\frac12a(t_0)^2,
  \qquad t_0\leq t<T_*.
\]
The integral before \(t_0\) is finite by classical smoothness, and the display
pays the rest of \(D_*\).  Equation
\eqref{eq:BodyEnstrophyAndHtwoFromDissipation} then gives a uniform \(H^1\)
roof, hence a uniform \(L^3\) roof, through the alleged last time.  This
contradicts the endpoint escape already used in
\eqref{eq:BodyFiniteEndpointForcesCriticalEscape}.  Thus no
\(T<T_*\) can exceed the displayed horizon.
\end{proof}

For a general datum, the energy balance stops one row below \(D_*\).  It pays
only
\[
  \int_0^{T_*}Z(t)\,dt<\infty.
\]
Interpolating \(D\leq Z^{1/2}Y^{1/2}\) inside
\eqref{eq:BodyFirstRectangleStretchingEstimate} yields
\begin{equation}\label{eq:BodyCubicEnstrophyObstruction}
  Z'(t)+\nu Y(t)
  \leq C\nu^{-3}Z(t)^3.
\end{equation}
The time integral of \(Z\) does not control either \(\int Z^3\) or
\(\sup Z\).  The energy and critical calculations have met at the same
boundary without crossing it.

They still constrain the shape of any failure.  Interpolation between energy
and enstrophy gives
\[
  H(t)^2
  =\frac14\norm{\Lambda^{1/2}u(t)}_2^4
  \leq\frac14\norm{u(t)}_2^2 Z(t),
\]
and a second time integration of the energy row gives
\begin{equation}\label{eq:BodyCriticalHeightTimeOccupancy}
  \int_0^{T_*}H(t)^2\,dt
  \leq\frac{\norm{u_0}_2^4}{8\nu}.
\end{equation}
Large critical height may occur, but it cannot occupy much time.  At the same
instant, interpolation between the energy height and the critical dissipation,
together with the energy ceiling in
\Cref{prop:ClassicalKineticEnergyIdentity}, gives, for a nonzero datum,
\begin{equation}\label{eq:BodyCriticalDissipationLowerFromHeight}
  \norm{\Lambda^{1/2}u(t)}_2
  \leq
  \norm{u(t)}_2^{2/3}
  \norm{\Lambda^{3/2}u(t)}_2^{1/3}
  \quad\Longrightarrow\quad
  D(t)\geq\frac{8H(t)^3}{\norm{u_0}_2^4}.
\end{equation}

Let the actual history climb from a critical height \(M\geq H(0)\) to
\(2M\).  Let \(t_1\) be its first arrival at \(2M\), and let \(t_0<t_1\) be
its last visit to \(M\) before that arrival.  On
\(I_M=(t_0,t_1)\), the height stays between \(M\) and \(2M\).  The occupancy
bound derived from \Cref{prop:ClassicalKineticEnergyIdentity} gives
\begin{equation}\label{eq:BodyCriticalDoublingIntervalLength}
  |I_M|
  \leq\frac{\norm{u_0}_2^4}{8\nu M^2}.
\end{equation}
Meanwhile
\(|\mathcal P_{1/2}|\leq C\sqrt{2H}\,D\).  Integrating the critical balance
over this first-passage interval gives
\begin{equation}\label{eq:BodyCriticalDoublingDissipationCost}
  M
  =\int_{I_M}H'(t)\,dt
  \leq2C\sqrt M\int_{I_M}D(t)\,dt,
  \qquad
  \int_{I_M}D(t)\,dt\geq\frac{\sqrt M}{2C}.
\end{equation}
This is the rigorous Zeno remnant.  Higher critical records are crossed in
shorter intervals, while every completed doubling carries a larger lower cost
in critical dissipation.  That lower cost exposes the pressure on \(D_*\); it
does not provide the missing upper budget.

Nor can an unpaid budget remain inside a fixed range of frequencies.  For
\(K\geq1\), set
\[
  D_{\leq K}(t)
  :=\int_{|\xi|\leq K}|\xi|^3|\widehat u(t,\xi)|^2\,d\xi.
\]
Since \(|\xi|^3\leq K|\xi|^2\) on this region, the kinetic-energy row pays
\begin{equation}\label{eq:BodyFixedFrequencyCriticalDissipation}
  \int_0^{T_*}D_{\leq K}(t)\,dt
  \leq K\int_0^{T_*}Z(t)\,dt
  \leq\frac{K\norm{u_0}_2^2}{2\nu}<\infty.
\end{equation}
If \(D_*\) diverges, its unpaid part moves beyond every fixed \(K\).  The
lowest endpoint obstruction is ultraviolet: it belongs to the same narrowing
history with which the vortex argument began.

The equivalence
\eqref{eq:BodyFirstRectangleCriticalDissipationEquivalence} concerns only the
lowest rectangle.  It does not promote that rectangle through every temporal
and spatial order.  To carry the complete pressure-coupled Tower to one common
terminal datum, the argument below needs the corresponding cutoff-independent
control at each finite coordinate.  We state that exact remaining input rather
than treating algebraic exposure as its proof.

\begin{assumption}[Whole-terminal rectangle condition]\label[assumption]{ass:BodyWholeTerminalRectangleCondition}
Let \(u_0\in\Ssigma\), let \(\nu>0\), and let \((u,p)\) be the maximal
classical solution generated by \((u_0,\nu)\) on \([0,T_*)\).  If
\(T_*<\infty\), then for every finite \(N,M\in\Nzero\),
\begin{equation}\label{eq:BodyWholeTerminalRectangleEstimate}
  \sup_{0<T<T_*}\mathfrak R_{N,M}(u;T)<\infty.
\end{equation}
The pressure rows are generated by
\eqref{eq:BodyTemporalPressureRecurrence} with the fixed decaying
normalization.
\end{assumption}
Here \((N,M)\) is fixed before the supremum over \(T<T_*\) is taken.  The bound
may depend on that finite pair.  Every rectangle still belongs to the same
history, and a larger rectangle must return the identical derivatives on an
overlap.

Assume now that the maximal history has a finite endpoint and that
\Cref{ass:BodyWholeTerminalRectangleCondition} holds on this branch.  Fix one
finite pair \((N,M)\).  Every term in
\eqref{eq:BodyRectangleNorm} is the integral of a nonnegative squared norm, and
the intervals \((0,T)\) exhaust \(I_*\).  Thus, for any one integrand \(f\),
\[
  \int_{I_*}f(t)\,dt
  =
  \lim_{T\uparrow T_*}\int_0^T f(t)\,dt.
\]
The cutoff-independent bound keeps this limit finite for every adjacent row in
the chosen rectangle:
\begin{equation}\label{eq:BodyWholeTerminalAdjacentRows}
  V_n\in L^2(I_*;H^{M+1}),
  \qquad
  V_{n+1}=\partial_tV_n\in L^2(I_*;H^{M-1}),
  \qquad 0\leq n\leq N.
\end{equation}
The pressure reconstruction in
\Cref{thm:WholeTerminalCompatibleRectangleEstimate} uses the Riesz bounds and
finite-order Sobolev products to recover the matching rows
\begin{equation}\label{eq:BodyWholeTerminalPressureRows}
  Q_n\in L^1(I_*;H^M),
  \qquad 0\leq n\leq N.
\end{equation}
The velocity and normalized pressure have now reached the whole interval as
one finite pressure-complete rectangle \(\cR_{N,M}[u_0]\).

Enlarge the rectangle.  The fluid does not acquire a second velocity, and a
derivative already reached in the smaller rectangle does not change its value.
If \(N'\leq N\) and \(M'\leq M\), let
\(\rho_{N',M'}^{N,M}\) forget the extra rows and read the retained rows at the
lower Sobolev height.  Every retained entry is still the same derivative of
the same \(u\), so
\begin{equation}\label{eq:BodyRectangleRestrictionCompatibility}
  \rho_{N',M'}^{N,M}\cR_{N,M}[u_0]
  =\cR_{N',M'}[u_0].
\end{equation}
This equality is the compatibility.  It lets every finite view be joined to
every larger view without gluing different histories:
\begin{equation}\label{eq:DatumSpecificCompatibleIntersection}
  \cI[u_0]
  :=\bigl(\cR_{N,M}[u_0]\bigr)_{N,M\in\Nzero},
  \qquad
  \rho_{N',M'}^{N,M}\cR_{N,M}[u_0]
  =\cR_{N',M'}[u_0].
\end{equation}
The notation creates no new field.  Given a finite temporal order and spatial
height, choose one rectangle large enough to see that derivative.  Every larger
rectangle returns the same derivative of the same \(u\).

\begin{proposition}[The mixed intersection]\label[proposition]{prop:BodyCompatibleProjectiveAssembly}
Under \Cref{ass:BodyWholeTerminalRectangleCondition}, the velocity component
of \(\cI[u_0]\) satisfies
\begin{equation}\label{eq:BodyProducedMixedIntersection}
  u\in
  \bigcap_{r=0}^{\infty}
  \bigcap_{m=0}^{\infty}
  H^r(I_*;H^m(\R^3)).
\end{equation}
In particular,
\begin{equation}\label{eq:CompleteSpatialRow}
  u\in\bigcap_{m=0}^{\infty}L^2(I_*;H^m(\R^3)).
\end{equation}
\end{proposition}

\begin{proof}
Fix finite \(r,m\), and choose \(N\geq r\) and \(M\geq m\).
The upper edge of this one rectangle gives
\[
  \partial_t^nu\in L^2(I_*;H^{M+1})
  \hookrightarrow L^2(I_*;H^m)
  \qquad(0\leq n\leq r).
\]
These are the distributional time derivatives of the same velocity field, so
\(u\in H^r(I_*;H^m)\).  Since \(r\) and \(m\) were arbitrary, every finite
coordinate belongs to the intersection; \(r=0\) gives
\eqref{eq:CompleteSpatialRow}.
\end{proof}
